\documentclass[12pt]{article}

\usepackage[margin=1.12in]{geometry}
\usepackage{amsmath,amssymb,amsthm,mathtools}
\usepackage{booktabs}
\usepackage{hyperref}

\newtheorem{proposition}{Proposition}
\newtheorem{theorem}{Theorem}
\theoremstyle{remark}
\newtheorem{remark}{Remark}

\newcommand{\C}{\mathbb C}
\newcommand{\eps}{\varepsilon}
\newcommand{\rank}{\operatorname{rank}}
\newcommand{\dist}{\operatorname{dist}}
\newcommand{\linf}{\ell^\infty}

\setlength{\emergencystretch}{3em}
\hypersetup{
  pdftitle={Quantitative Readout Accountability for Finite Centers},
  pdfsubject={Finite-center readout stability, sufficiency, binary-rank obstruction},
  pdfkeywords={finite centers, readout accountability, condition number, Le Cam deficiency, Blackwell comparison, binary rank}
}

\title{Quantitative Readout Accountability\\
for Finite Centers}
\author{}
\date{Live review note of June 20, 2026}

\begin{document}
\maketitle

\begin{abstract}
The exact Team 1 support/readout criterion says that finite central rank is
recoverable only when the supplied readouts are faithful on the physical
finite-center contrast space, after null/gauge/inadmissible directions have
been quotiented.  This note records the finite-precision strengthening.  If
\[
  R:V_F=\ell^\infty(F)/\mathbb C1\longrightarrow W
\]
is the declared central readout map, then exact accountability is
\(\ker R=0\).  Noisy accountability additionally requires a lower inverse
modulus
\[
  \alpha(R)=\inf_{\|v\|=1}\|Rv\|>0
\]
large enough for the error budget.  If two contrasts are compatible with the
same data ball of radius \(\varepsilon\), then their distance is at most
\(2\varepsilon/\alpha(R)\).  Thus a formally separating but badly conditioned
readout is not a referee-grade finite-rank or binary-rank measurement unless
its separation margin dominates the stated uncertainty.  This is standard
finite-dimensional inverse stability; the possible value is the explicit
black-hole/AQFT finite-center warning.
\end{abstract}

\paragraph{Status.}
This is non-frozen external-review metadata for Team 1.  It does not modify
any frozen manuscript, Lean archive, public mirror manifest, or SHA-256
package.

\section{Primary sources}

The boundary below was checked against Blackwell's comparison of experiments
\cite{Blackwell1951,Blackwell1953}, Le Cam's sufficiency and approximate
sufficiency \cite{LeCam1964}, Petz sufficiency for von Neumann algebras
\cite{Petz1986}, and Buscemi's quantum statistical comparison theorem
\cite{Buscemi2012}.  These are positive sufficiency/deficiency frameworks.
No novelty is claimed for the linear stability estimate.

\section{Exact and noisy readout}

Let \(F\) be finite and set
\[
  V_F=\ell^\infty(F)/\mathbb C1 .
\]
Let \(R:V_F\to W\) be the finite-dimensional readout map generated by the
declared central probes, effects, sector records, recovery task, spectroscopy
channel, or other finite-center observation.  Fix norms on \(V_F\) and \(W\).
Define
\[
  \alpha(R)=\inf_{\|v\|=1}\|Rv\|.
\]

\begin{proposition}[exact accountability]
\label{prop:exact}
The readout \(R\) separates all finite central contrasts if and only if
\(\alpha(R)>0\).  In finite dimension this is equivalent to \(\ker R=0\), and
hence to \(\rank R=\dim V_F\).
\end{proposition}

\begin{proof}
If \(\alpha(R)>0\), then \(Rv=0\) implies
\(\alpha(R)\|v\|\le\|Rv\|=0\), so \(v=0\).  Conversely, if \(\ker R=0\), the
continuous function \(v\mapsto\|Rv\|\) has a positive minimum on the compact
unit sphere of \(V_F\).
\end{proof}

\begin{proposition}[noisy accountability]
\label{prop:noisy}
Assume \(\alpha(R)>0\).  If two contrasts \(v,v'\in V_F\) are compatible with
the same observed readout \(y\in W\) within error \(\varepsilon\),
\[
  \|Rv-y\|\le\varepsilon,\qquad
  \|Rv'-y\|\le\varepsilon,
\]
then
\[
  \|v-v'\|\le {2\varepsilon\over\alpha(R)}.
\]
Consequently contrasts separated by distance at least \(\delta\) are
distinguishable at error \(\varepsilon\) only if
\[
  \alpha(R)\delta>2\varepsilon .
\]
\end{proposition}

\begin{proof}
The triangle inequality gives
\[
  \|R(v-v')\|
  \le \|Rv-y\|+\|Rv'-y\|
  \le 2\varepsilon .
\]
By definition of \(\alpha(R)\),
\[
  \alpha(R)\|v-v'\|\le\|R(v-v')\|,
\]
which proves the estimate.  The distinguishability condition is the same
inequality read contrapositively for a required contrast resolution \(\delta\).
\end{proof}

\begin{proposition}[null or gauge quotient]
Let \(N\subset V_F\) be a declared null, gauge, or inadmissible subspace and
let \(\bar R:V_F/N\to W\) be the induced readout.  Propositions
\ref{prop:exact} and \ref{prop:noisy} hold with \(V_F\) replaced by \(V_F/N\)
and \(\alpha(R)\) replaced by \(\alpha(\bar R)\).
\end{proposition}

\begin{proof}
Apply the previous two propositions to the finite-dimensional quotient space.
\end{proof}

\section{Binary-rank consequence}

A noisy operational claim of finite central rank, and especially a claim of
binary rank \(q=2\), requires more than a formally separating family of
questions.  It requires:
\[
\begin{array}{p{0.34\linewidth}p{0.52\linewidth}}
\toprule
Premise & Role \\
\midrule
Physical contrast space \(V_F/N\) & States which finite central directions are
physical rather than null, gauge, or inadmissible. \\
\addlinespace
Full exact rank & Ensures no physical finite central contrast is in the
kernel. \\
\addlinespace
Lower inverse modulus \(\alpha(\bar R)\) & Converts formal separation into a
stable error-budget statement. \\
\addlinespace
Resolution inequality \(\alpha(\bar R)\delta>2\varepsilon\) & Ensures the
claimed rank-breaking contrast is larger than uncertainty. \\
\addlinespace
Maximality, finite-alphabet law, or direct spectroscopy & Prevents a stable
binary readout from being merely a binary quotient of a larger hidden center.
\\
\bottomrule
\end{array}
\]

Thus exact accountability answers the algebraic question, while quantitative
accountability answers the referee's finite-precision question.

\section{Referee question}

\begin{quote}
In the target black-hole/AQFT reconstruction setting, what is the lower
readout modulus or equivalent Le Cam/deficiency/error-budget statement that
makes the finite central rank claim stable?
\end{quote}

If no such modulus or deficiency bound is supplied, Team 1 should treat the
rank or \(q=2\) statement as conditional.  If the bound is supplied, the
finite-center obstruction has an explicit positive exit at the stated
precision.

\begin{thebibliography}{9}

\bibitem{Blackwell1951}
D. Blackwell,
\emph{Comparison of Experiments},
Proceedings of the Second Berkeley Symposium on Mathematical Statistics and
Probability, 1951, pp. 93--102,
\href{https://projecteuclid.org/euclid.bsmsp/1200500222}{Project Euclid}.

\bibitem{Blackwell1953}
D. Blackwell,
\emph{Equivalent Comparisons of Experiments},
Ann. Math. Statist. 24 (1953), 265--272.

\bibitem{LeCam1964}
L. Le Cam,
\emph{Sufficiency and Approximate Sufficiency},
Ann. Math. Statist. 35 (1964), 1419--1455.

\bibitem{Petz1986}
D. Petz,
\emph{Sufficient Subalgebras and the Relative Entropy of States of a von
Neumann Algebra},
Commun. Math. Phys. 105 (1986), 123--131,
\href{https://projecteuclid.org/journals/communications-in-mathematical-physics/volume-105/issue-1/Sufficient-subalgebras-and-the-relative-entropy-of-states/cmp/1104115260.pdf}{Project Euclid PDF}.

\bibitem{Buscemi2012}
F. Buscemi,
\emph{Comparison of Quantum Statistical Models: Equivalent Conditions for
Sufficiency},
Commun. Math. Phys. 310 (2012), 625--647,
\href{https://arxiv.org/abs/1004.3794}{arXiv:1004.3794}.

\end{thebibliography}

\end{document}
